\section{Evidence and Evaluation Framework}
\label{sec:framework}

\subsection{Evidence source: LongMemEval}
\label{sec:longmemeval}

Rather than generating synthetic facts (as in our preliminary experiments),
we use evidence from the LongMemEval benchmark \citep{wu2024longmemeval}. LongMemEval
provides 500 question-answer pairs across 6 categories, each grounded in
realistic multi-session conversations between a user and an LLM assistant.

Each fact consists of:
\begin{itemize}
  \item A \textbf{conversation excerpt} (the ``evidence'') containing the information
  \item A \textbf{question} that can only be answered from the evidence
  \item An \textbf{expected answer} with keywords for automated verification
  \item A \textbf{category} (knowledge-update, single-session-user, single-session-assistant,
    single-session-preference, temporal-reasoning, multi-session)
\end{itemize}

This provides ecologically valid facts --- they are naturally embedded in
conversation, not injected as artificial needles.

\subsection{Factorial evidence design (2$\times$2)}
\label{sec:factorial}

We vary two dimensions of evidence preparation:

\textbf{Filtering} (complete vs filtered):
\begin{itemize}
  \item \textit{Complete}: all 500 questions included
  \item \textit{Filtered}: only 234 questions from answerable categories (single-session-user,
    single-session-assistant, knowledge-update, single-session-preference).
    Excludes temporal-reasoning and multi-session, which require capabilities
    beyond single-context recall.
\end{itemize}

\textbf{Truncation} (full vs chopped):
\begin{itemize}
  \item \textit{Full}: evidence messages preserved in their entirety
  \item \textit{Chopped}: messages truncated to realistic lengths (simulating context limits
    in real systems)
\end{itemize}

This yields four experimental modes:

\begin{table}[h]
\centering
\caption{Four experimental modes from the 2$\times$2 factorial design.}
\label{tab:modes}
\begin{tabular}{lllcc}
\toprule
Mode & Filtering & Truncation & Questions & Max density (facts/kTok) \\
\midrule
R1 & Complete & Full    & 500 & 0.10 \\
R2 & Complete & Chopped & 500 & 0.42 \\
R3 & Filtered & Full    & 234 & 0.16 \\
R4 & Filtered & Chopped & 234 & 0.79 \\
\bottomrule
\end{tabular}
\end{table}

\subsection{Context construction}
\label{sec:context-construction}

For a given density $\delta$, we embed $N = \delta \times L$ evidence items into a context of
$L$ kilotokens (190 kTok for calibration and single-pass experiments). Evidence
items are placed at uniformly distributed positions. The remaining space is
filled with realistic padding --- real LLM conversation sessions from the
LongMemEval pool (18{,}255 sessions available) --- maintaining ecological
validity rather than using synthetic filler.

Context construction is deterministic (seed=42), ensuring reproducibility
across runs. The padding pool provides $\sim$46M tokens of real conversations.

\subsection{Density sweep}
\label{sec:density-sweep}

We define \textbf{fact density} $\delta = N/L$, where $N$ is the number of embedded facts
and $L$ the context size in kilotokens (kTok). Rather than testing a single
density, we sweep across values:

\begin{table}[h]
\centering
\caption{Density sweep used in the calibration and single-pass experiments.}
\label{tab:density-sweep}
\begin{tabular}{ccc}
\toprule
$N$ (facts) & Context & $\delta$ (facts/kTok) \\
\midrule
4   & 190 kTok & 0.02 \\
10  & 190 kTok & 0.05 \\
20  & 190 kTok & 0.11 \\
40  & 190 kTok & 0.21 \\
60  & 190 kTok & 0.32 \\
80  & 190 kTok & 0.42 \\
150 & 190 kTok & 0.79 \\
\bottomrule
\end{tabular}
\end{table}

This reveals saturation effects: at what density does the model's recall
capacity plateau? The scripts use shorthand \code{dN} (e.g., \code{d80} = 80 facts in
190 kTok = $\delta$ 0.42); in the text, we report density as $\delta$.

\subsection{Evaluation protocol}
\label{sec:eval-protocol}

\textbf{Q\&A phase}: For each fact, the LLM (Claude Haiku 4.5) is presented with
the full context as conversation history and asked to recall the specific
information. Multiple questions can be asked in a single prompt; we denote
$Q$ the number of \textbf{questions per prompt} and test $Q=1$, $Q=5$, and $Q=10$.

\textbf{Judge phase}: A separate LLM judge compares the model's answer against the
expected answer keywords. Binary recall (correct/incorrect) and accuracy
(correct and precisely matching) are computed.

\textbf{Grep validation}: As a free upper bound, we check whether fact keywords
appear verbatim in the context. If grep doesn't find them, the LLM cannot
possibly recall them. The gap between grep recall and LLM recall quantifies
the ``present but ignored'' phenomenon.

Both Q\&A and judge phases use Claude Haiku 4.5 via the Anthropic Batch API.

\subsection{Repeatability}
\label{sec:repeatability}

All key results are validated with 3 independent runs. We report mean $\pm \sigma$.
For the single-pass compaction experiment (\S\ref{sec:singlepass}), 3 runs at $\delta=0.42$/$Q=5$ yield $\sigma \leq 2.6$pp
across all compaction levels, confirming that observed effects (7--70pp) are
far larger than measurement noise.
